\section{Interfaces and Downstream Integration}
\label{sec:integration}

\plain{In brief: how the pipeline is invoked, how it embeds in other software, and what a downstream tool is told about the geometry it receives.}

\paragraph{Remote-volume workflow.} The public \texttt{scrollunwrap} entry point streams an OME-Zarr region from anonymous or authenticated S3 storage, sends the selected voxels to the in-memory C mesher, and can continue through preview and \texttt{tifxyz} output without first materializing a TIFF cube grid. It also supports local arrays, mesh-only operation, independent cube and flattener timeouts, and graceful cube skipping. All outputs carry the true $(x,y,z)$ origin and scale, so the resulting OBJ is immediately viewable and its atlas indexes the source volume rather than an internal cube convention.

\paragraph{Unified quad-ribbon command.} The whole-sheet lane is exposed as \texttt{quadribbon <mesh\_dir> <out\_dir> [config.json]}. It recursively accepts the world-frame per-cube mesh pile and applies the fixed production order of fit, metric solve, untangle, metric re-solve, strip bake, and composite. Stage artifacts are resumable, and the production constants --- axis, pitch, CT source, bake window, dark threshold --- come from one configuration file rather than from tuning flags on the command line, so that a reported number belongs to a stated configuration rather than to a session.

\paragraph{Embeddable implementation.} One versioned public C header exposes cleanup, split, topology, reconstruction, remeshing, weld, and full volume-to-mesh operations. A consumer can use a static library, a shared library, or a runtime-loaded ABI table; cross-runtime allocation, structured status, progress, timeout, and cancellation are all explicit. The CMake build is exercised both as an embedded consumer and on Linux/GCC. The MLS collapse can run through the reference C path, OpenMP, an optional FP32 CUDA backend, or a research CubeCL CUDA backend; every backend applies the same guards and is required to produce the same geometry, so a backend may be faster but never more permissive.

\paragraph{Application handoff.} A clean local component enters \texttt{villa-patch}: ScrollFiesta preserves its intermediate ribbon, quilting, and snap artifacts, and Villa performs the authoritative final flatten. Larger uncertain surfaces enter VC3D through \texttt{ScrollFiesta: Generate Spiral Hints}, which creates a new \texttt{scroll-hint}/\texttt{unverified} segment, supports progress and cancellation, reloads the surface, and can add it to the active spiral fit; it never overwrites its source. The separation matters because the two consumers want different things: a patch tool wants geometry whose wrap identity is settled, and a whole-scroll fitter wants candidate surfaces whose identity is still open.

\paragraph{Output labels.} Geometry and textures distinguish released prediction, direct fit, harmonic or cylindrical fill, moved-by-CT vertices, suspended order relations, primary and secondary ply, bounded raster void-fill, and masked output. A consumer receiving a partial result can therefore use it without confusing a generated continuation for observed papyrus, and a consumer that disagrees with one of our thresholds can re-derive its own conclusion from the labels rather than from the picture.
